\scanpage{18}
\mspage{17}

Il n'y a pas à y revenir. Je vais plutôt m'attacher à présent aux interprétations en termes d'ombres (grandes ou petites). Il est entendu que dans un premier temps, je m'attache aux \unsure{aspects|aspect} ``algébriques'' ou ``triviaux'' d'un formalisme des figures, et me contente de dégager \unsure{les|le} \unsure{bons|bon} \unsure{axiomes|axiome}, de telle façon que

1\textsuperscript{o}) Ceux-ci \unsure{\deleted{doivent}|doivent} valables dans tout ``atelier'' de la forme $\mathrm{Fig}(\deleted{\unsure{E|L}}L)$, formé des figures ensemblistes dans un ensemble donné $L$ (sans conditions de finitude d'aucune sorte),

et

2\textsuperscript{o}) qu'ils soient applicables aussi, disons, aux figures \uline{simpliciales} affines finies, dans un espace affine sur $\mathbb{R}$ (ou, plus gén\textsuperscript{t}., sur un corps ordonné quelconque).

Pour expliciter, en termes de l'ens ordonné $(\unsure{M|\mathcal{M}},\leq)$, \unsure{les ens|l'ens} $\unsure{\mathcal{F}|\mathfrak{F}|\mathscr{F}}$ des figures, en tant qu'ens de parties de $M$, \deleted{il \unsure{est}} ou du moins des \uline{figures de type fini}
\[
\begin{aligned}
 F\in\mathcal{F}\ \text{de t.f} \quad&\Longleftrightarrow\quad
 \widetilde F=\deleted{\unsure{\text{fibres}|\text{fibre}}}\,\text{réunion finie}\\
 &\hspace{4.9em}\text{d'un}\,M_{\leq \unsure{x_i|X_i}}\\
 &\hspace{3.0em}\text{(pour un nb \uline{fini}}\\
 &\hspace{3.5em}\text{de strates de }\deleted{\unsure{M|x_i|X}}\text{)}
\end{aligned}
\tag*{(1.34)}
\]
